\section{Conclusion}
The conclusion summarizes the results of the research, evaluates the achieved results and formulates the prospects for further work. In the course of the study, the methodology of agent-based modeling was developed to analyze the behavior of the financial system under market shocks. The results obtained showed that:

\begin{enumerate}
    \item Agent-based modeling effectively reproduces the dynamics of financial markets, demonstrating characteristic patterns under abrupt changes in conditions.
    \item The introduction of shock events leads to significant changes in volatility and liquidity, confirming hypotheses derived from theoretical studies (Danielsson, 2005; Brunnermeier, 2009).
    \item The parametric analysis allowed us to identify the key factors that influence the stability of the system, as well as to determine the scope of application of the model for the development of practical strategies for market stabilization.
\end{enumerate}

The main limitations of the study are related to the simplified model of the market system, which requires further expansion of the simulator's functionality and deeper empirical substantiation. Prospects for further research include the integration of additional factors (e.g., the impact of regulatory measures, intermarket interactions) and the application of the model to the analysis of real market data.